% --- DIAPOSITIVA 2A: INTRODUCCIÓN - CONTEXTO Y MOTIVACIÓN (I) ---
\begin{frame}{Introducción: Contexto y Motivación (I)}
	\begin{block}{Polinomios Ortogonales Clásicos}
		\begin{itemize}
			\item Herramienta fundamental en matemáticas, física e ingeniería.
			\item Teoría bien establecida: recurrencia de tres términos, ceros reales y entrelazados.
		\end{itemize}
	\end{block}
	\pause
	\begin{alertblock}{Generalización: Polinomios Ortogonales de Sobolev} % Tamaño de fuente reducido
		\begin{itemize}
			\item Productos internos que consideran derivadas de los polinomios.
			\item Motivación: Mejorar aproximación de funciones y sus derivadas.
			\item Fenómenos nuevos:
			      \begin{itemize}
				      \item Ceros pueden ser complejos y ``errantes''.
				      \item El operador de multiplicación por $z$ no es autoadjunto; su norma (y la de sus secciones de Hessenberg, vía valores singulares\cite{singularValuesOfHessenbergMatrices}) controla la acotación de los ceros.
			      \end{itemize}
		\end{itemize}
	\end{alertblock}
\end{frame}

% --- DIAPOSITIVA 2B: INTRODUCCIÓN - CONTEXTO Y MOTIVACIÓN (II) ---
\begin{frame}{Introducción: Contexto y Motivación (II)}
	\begin{block}{Importancia del Estudio de los Ceros de Polinomios de Sobolev} \footnotesize % Tamaño de fuente reducido
		\begin{itemize}
			\item Diseño de reglas de cuadratura (fórmulas de integración numérica) de Sobolev estables.
			\item Estabilidad de métodos numéricos para Ecuaciones en Derivadas Parciales (EDP).
			\item Conexiones con teoría espectral (estudio de autovalores) y teoría del potencial.
			\item Comprensión de fenómenos de transición de fase en matrices aleatorias con aplicaciones (física estadística, teoría cuántica).
		\end{itemize}
	\end{block}
\end{frame}

% --- DIAPOSITIVA 3: INTRODUCCIÓN - OBJETIVOS (I) --- % Título modificado
\begin{frame}{Introducción: Objetivos del TFG (I)} % Título modificado
	\footnotesize % Tamaño de fuente reducido para toda la diapositiva
	Este trabajo se centra en:
	\begin{enumerate}
		\item \textbf{Recopilar y estudiar la teoría} de polinomios ortogonales de Sobolev, unificando conceptos sobre la acotación de sus ceros y el comportamiento espectral de las matrices asociadas.
		      \pause
		\item \textbf{Diseñar algoritmos} para la visualización del conjunto de ceros a partir de las representaciones matriciales del operador de multiplicación por la variable.
		      \pause
		\item \textbf{Implementar computacionalmente} en \texttt{Maple} dichos algoritmos, automatizando la construcción de matrices, el cálculo de parámetros espectrales relevantes y la visualización interactiva.
		      \pause
		\item \textbf{Analizar y verificar resultados teóricos} mediante experimentación numérica, contrastar cotas conocidas y explorar nuevos patrones para formular conjeturas.
	\end{enumerate}
\end{frame}

% --- NUEVA DIAPOSITIVA 3B: INTRODUCCIÓN - OBJETIVOS (II) - ENFOQUE Y ALCANCE ---
\begin{frame}{Introducción: Objetivos del TFG (II) - Enfoque y Alcance}
	\footnotesize % Tamaño de fuente reducido para toda la diapositiva
	La investigación se desarrolla en dos \textbf{frentes complementarios}:
	\begin{itemize}
		\item \textit{Resultados analíticos:} Avances en teoría espectral, localización de ceros y condiciones de acotación del operador de multiplicación por la variable \cite{lagomasino2001sobolev, DuranSaff2001, escribano2025zeros}.
		\item \textit{Experimentación numérica y visualización:} Uso de \texttt{Maple} para generar matrices, analizar su espectro, visualizar ceros y validar conjeturas (e.g., convergencia del segundo valor singular más grande).
	\end{itemize}
	\pause
	Este TFG desarrolla un \textbf{marco teórico-computacional unificado} que:
	\begin{itemize}
		\item Formaliza la construcción de matrices (de momentos, del operador de multiplicación) para productos de Sobolev de primer orden.
		\item Relaciona la acotación de los ceros con la norma y ciertos valores singulares (el primero y el segundo más grandes) de las matrices del operador.
		\item Implementa algoritmos en \texttt{Maple} para cálculo y visualización 2D de ceros, soportes de las medidas y espectro.
		\item Presenta ejemplos (casos reales, circunferencias, mezclas con términos discretos) ilustrando la influencia de la geometría de los soportes.
	\end{itemize}
	\pause
	\vspace{0.2cm}
	\begin{center}
		\textit{Se persigue clarificar la interacción entre geometría del soporte, álgebra matricial y análisis funcional, y proponer nuevas conjeturas sobre comportamientos asintóticos (radios espectrales, segundos valores singulares).}
	\end{center}
\end{frame}